\subsection{Revised candidate prior and preparation interface}
\label{sec:v2-method}

The original regime rule is retained as a legacy control. A Roofline bound characterizes a resource limit; it does not rank fused implementations, account for launch overhead, or guarantee that removing a materialized tensor improves latency. The revised, opt-in SCIC-v2 uses the analytical quantities as candidate-cost features rather than excluding a candidate because a shape lies in a particular intensity band. Memory-classified shapes are eligible for calibration.

The bounded implementation pool contains five schedules: three Triton projection schedules with tiles $(16,64,64)$ and split factors 1, 4, and 8, and two cuBLAS projections followed by a Triton residual/RMSNorm kernel using four or eight warps. Split partials accumulate in FP32. A second kernel iterates over partial buffers, rounds the projection to FP16, and applies the specified residual convention. The partial buffers incur additional traffic; these schedules are not assumed to reduce DRAM traffic. The HF convention rounds the residual sum before normalization and casts normalized values before the affine multiplication. The vLLM convention computes statistics and affine multiplication from the FP32 residual sum, then stores the FP16 residual. These conventions have separate validation and training domains.

For each candidate, ridge regression predicts log synchronized host-batch latency, with regularization $10^{-3}$. Shape, execution mode, weight reuse, and schedule features form the base model. Complete-chain FLOPs, explicit candidate traffic, analytical Roofline time, and intensity relative to the ridge supplement it. Feature scaling and fitting use training geometries only. The 95th percentile absolute log residual under leave-one-training-geometry-out fitting provides an empirical uncertainty margin; this margin is not a distribution-free confidence guarantee. The no-AI ablation removes the analytical features while preserving training data, fitting, and probe rules.

Preparation starts with the training-selected fixed incumbent and the predicted challenger. It expands the pool if an unexamined candidate's lower uncertainty bound can improve the measured winner by at least 2\%. Candidates must pass numerical checks before being probed. A candidate replaces the incumbent only when projected savings over the declared reuse horizon amortize incremental preparation. Compilation, numerical validation, and probing are recorded separately; probe-only reductions cannot establish a reduction in total preparation cost.

The \texttt{prepare} interface returns a frozen callable before outer compilation or CUDA Graph capture. Its cache signature includes shape and stride, precision settings, device, software and implementation versions, normalization convention and epsilon, and weight identity/version. Inference tensors do not expose an automatic mutation counter; a caller changing frozen projection or normalization weights must bump the explicit weight-version argument and prepare again. The tested model adapters assume immutable loaded weights. Unsupported layouts use a recorded native fallback. The HF endpoint is traced into the outer graph through \texttt{triton\_op} and \texttt{wrap\_triton}; it has no per-layer exact-key dispatch or nested compiled projection call during steady execution. A vLLM worker prepares loaded weights after the framework's \texttt{load\_model} returns and before profiling/compilation/capture. Its serving operator declares residual mutation and updates that buffer without the legacy adapter's extra residual clone. Deployment remains limited to FP16, tensor-parallel size one, and the tested Qwen3 interfaces.

If serving preparation cannot admit a common candidate for the declared mixture, it restores that layer's original vLLM projection and forward method before tracing. The preparation report marks native deployment and its rejection reason; no runtime branch selects this fallback.
